\documentclass[11pt]{article}

\usepackage[T1]{fontenc}
\usepackage{lmodern}
\usepackage{amsmath,amssymb,amsthm}
\usepackage[margin=1in]{geometry}
\usepackage{enumitem}
\usepackage{xcolor}
\usepackage{microtype}
\usepackage[hidelinks]{hyperref}

\newcommand{\R}{\mathbb{R}}
\newcommand{\Om}{\Omega}
\newcommand{\norm}[1]{\left\lVert #1\right\rVert}

\title{Paper 3, Theorem 2.1: Statement and Interface Amendment}
\author{Internal research note for the \texttt{ShenWork} formalization}
\date{14 July 2026}

\begin{document}
\maketitle

\begin{abstract}
Theorem~2.1(1) of Chen--Ruau--Shen, Part~II, claims uniform persistence under
the sole condition $m\geq1$, but its proof treats only $a=b=0$ and $a,b>0$.
The omitted regime $a=0<b$ contains a spatially constant, positive, bounded,
global solution that decays to zero, so the printed statement is false.  The
correct parameter guard is the disjunction actually used in the proof.  A
second issue concerns the formal interface for part~(4): the paper's mass is
attached to the initial datum, whereas the first Lean version read the freely
stored value $u(0)$ of a positive-time trajectory.  This note gives the exact
counterexamples, the corrected statements, and the corresponding axiom-clean
Lean capstones.
\end{abstract}

\section{Source checked}

The source examined is the original submission, not a repository paraphrase:
\begin{quote}
Le Chen, Ian Ruau, and Wenxian Shen, \emph{Chemotaxis models with
signal-dependent sensitivity and a logistic-type source, II: Persistence and
stabilization}, arXiv:2604.02599v1, 3~April 2026.
\end{quote}
Theorem~2.1 is titled ``Uniform persistence.''  Section~4.1 proves part~(1),
and Section~4.4 proves part~(4).  The two corrections below have different
status: the first is a genuine parameter omission in the printed theorem; the
second repairs the representation of the paper's initial mass in the Lean
trajectory API.

\section{The omitted pure-decay regime in part (1)}

The printed part~(1) assumes only $m\geq1$ and asserts, for every globally
defined, bounded, positive solution,
\[
  \liminf_{t\to\infty}\inf_{x\in\Om}u(t,x)>0.
\]
Section~4.1 nevertheless divides its proof into exactly two cases:
\begin{enumerate}[label=\arabic*.]
\item $a=b=0$;
\item $a>0$ and $b>0$.
\end{enumerate}
It never treats the allowed parameter regime $a=0<b$, and in that regime the
conclusion is false.

\subsection{Concrete counterexample}

On any bounded Neumann domain, let $c>0$, $a=0$, and $b>0$, and define the
spatially constant functions
\[
  u(t,x)=\bigl(c^{-\alpha}+\alpha b t\bigr)^{-1/\alpha},
  \qquad
  v(t,x)=\frac{\nu}{\mu}u(t,x)^{\gamma}.
\]
All spatial derivatives vanish.  Hence the diffusion, chemotaxis-divergence,
and boundary terms vanish, the elliptic equation holds identically, and the
parabolic equation becomes
\[
  u_t=-b u^{1+\alpha}.
\]
The solution is classical, strictly positive at every finite positive time,
global, and bounded by $c$.  However,
\[
  \lim_{t\to\infty}\inf_{x\in\Om}u(t,x)=0.
\]
The Lean counterexample specializes to
$\alpha=\gamma=m=\mu=\nu=b=1$ and $a=\chi_0=0$, so that on positive time
\[
  u(t,x)=v(t,x)=\frac{1}{1+t}.
\]

\subsection{Corrected part (1)}

The paper-faithful correction is the following.

\begin{quote}
Assume $m\geq1$ and either $a=b=0$ or $a>0$, $b>0$.  Then every globally
defined, bounded, positive solution has a strictly positive asymptotic spatial
lower bound for $u$, and the displayed elliptic lower bound for $v$ follows.
\end{quote}

The remaining nonnegative-parameter regime $a>0$, $b=0$ does not produce a
decaying counterexample.  Integrating the equation and using both Neumann
conditions gives
\[
  M'(t)=aM(t), \qquad M(t):=\int_{\Om}u(t,x)\,dx.
\]
A positive solution therefore has exponentially growing mass and cannot be
globally bounded.  The persistence implication is vacuous in that regime.
For fidelity, the recommended statement records the two cases proved in
Section~4.1 instead of silently adding a vacuous third branch.

\section{The positive-time mass interface in part (4)}

Part~(4) is a minimal-model statement ($a=b=0$, $m=1$).  The paper starts from
an initial datum $u_0$, constructs its solution, and imposes the mass condition
on $u_0$.  This is unambiguous because $u(t)\to u_0$ as $t\to0^+$ and mass is
conserved.

The first Lean statement instead used
\[
  \texttt{HasInitialMass}(D,u,u_*)
  \quad:\Longleftrightarrow\quad
  \int_D u(0)=|D|u_*.
\]
But \texttt{IsPaper2GlobalClassicalSolution} constrains the trajectory only at
strictly positive times.  Its stored value $u(0)$ is free.  Replacing only that
zero-time slice preserves the full positive-time classical orbit, positivity,
and boundedness, while changing \texttt{HasInitialMass} arbitrarily.  Thus the
old Lean version of part~(4) is false for every admissible constants package.

The correct formal interface is the conserved physical mass
\[
  \forall t>0,\qquad \int_D u(t)=|D|u_*.
\]
Equivalently, one may quantify over $u_0$, require an actual initial-trace
relation, and impose the mass condition on $u_0$.  The positive-time version is
the shorter interface for the persistence proof and states exactly the
invariant it consumes.  This is a Lean-interface correction, not a claim that
the paper's mathematical part~(4) is false.

\section{Lean realization}

The concrete interval counterexample to part~(1) is
\begin{quote}
\texttt{ShenWork.Paper3.}\
\texttt{not\_Theorem\_2\_1\_part1\_intervalDomain\_pureDecay}
\end{quote}
in
\path{ShenWork/Paper3/IntervalDomainPersistencePart1StatementObstruction.lean}.
It constructs the genuine interval-domain trajectory, proves global classical
solvability and boundedness, computes the lower-envelope liminf as zero, and
refutes the unguarded statement.  The same file defines
\texttt{Theorem\_2\_1\_part1\_corrected} with the Section~4.1 parameter split.

For part~(4),
\begin{quote}
\texttt{ShenWork.Paper3.}\
\texttt{not\_intervalDomain\_Theorem\_2\_1\_part4\_anyConstants}
\end{quote}
in
\path{ShenWork/Paper3/IntervalDomainPersistenceMinimalMassInterfaceObstruction.lean}
proves failure of the zero-slice formulation for every
\texttt{Paper3Constants} package.  The corrected capstone is
\begin{quote}
\texttt{ShenWork.Paper3.}\
\texttt{Theorem\_2\_1\_part4\_intervalDomainM\_physicalMass\_proven}
\end{quote}
in
\path{ShenWork/Paper3/IntervalDomainPersistenceMinimalPhysicalMass.lean}.
It uses the paper-faithful $u^m$ domain, conserved positive-time mass, the
proved eventual upper bound, and the quantitative Neumann-resolver mass gap.

All three capstones print only the standard Lean/Mathlib axioms
\texttt{propext}, \texttt{Classical.choice}, and \texttt{Quot.sound}.

\section{Recommended amendment}

\begin{enumerate}[label=\arabic*.]
\item Amend Theorem~2.1(1) by adding
  $(a=b=0)\lor(a>0\ \text{and}\ b>0)$.
\item Retain the mathematical mass hypothesis of Theorem~2.1(4), but in Lean
  attach it to the initial datum plus initial trace, or use the equivalent
  conserved positive-time mass.
\item Do not use an axiom-clean proof as a substitute for semantic validation:
  both failures above are themselves provable without any nonstandard axiom.
\end{enumerate}

\end{document}
